\input format.tex
\usepackage{mdframed}
\usepackage{lipsum}

% all 4 borders
\newmdenv[linecolor=black70]{allfour}
% just top and bottom
\newmdenv[leftline=false,rightline=false]{topbot}
% just left and bottom
\newmdenv[topline=false,rightline=false]{leftbot}

\begin{document}

\setlength{\arrayrulewidth}{1pt}
\fontsize{9.3pt}{17pt}\selectfont
\color{gray2}

\vspace*{3mm}
\begin{center}
\fontsize{16pt}{\baselineskip}\selectfont\color{black70}\bf {一、肠道菌群基因检测结果}
\end{center}

\noindent{\fontsize{11pt}{11pt}\selectfont\color{black70}\bf {肠道菌群总体状况}}

\vspace*{3mm}
\noindent
\begin{minipage}[b]{.32\textwidth}
\fontsize{9.3pt}{10pt}\selectfont
\begin{tctabularx}{tabularx={p{4.6cm}}}
\parbox[c]{.9\hsize}{\centering\vskip4pt{\color{white}\bf 菌群多样性：
健康
} \vskip4pt}
\\\\\\\\\\\\\\
\\\\\\\\\\\\\\
\\\\\\\\
\parbox[c]{\hsize}{
\hspace*{1cm}{
\begin{picture}(1, 1)
\begin{tikzpicture}[bar interval width=0.8, bar width=0.70cm]
	\begin{scope}
	\clip[] (0cm, 0cm) rectangle (1.5cm, 5.7cm);
	\draw[fill=solomon0, solomon0] (0cm, 0 * 5.7cm) rectangle (1.5cm, 0.05 * 5.7cm) node {};
	\draw[fill=solomon1, solomon1] (0cm, 0.05 * 5.7cm) rectangle (1.5cm, 0.15cm * 5.7cm) node {};
	\draw[fill=solomon2, solomon2] (0cm, 0.15 * 5.7cm) rectangle (1.5cm, 0.3 * 5.7cm) node {};
	\draw[fill=solomon3, solomon3] (0cm, 0.3 * 5.7cm) rectangle (1.5cm, 1 * 5.7cm) node {};
	\end{scope}
	\draw (-0.35cm, 0cm) node[left, anchor = west, yshift=0mm] {\bahao\color{gray2} {0}};
	\draw (0.4cm, 0.025 * 5.7cm ) node[left, anchor = west, yshift=0mm] {\qihao\color{white}\bf {脆弱}};
	\draw (-0.35cm, 0.05 * 5.7cm) node[left, anchor = west, yshift=0mm] {\bahao\color{gray2} {5}};
	\draw (0.4cm, 0.1 * 5.7cm ) node[left, anchor = west, yshift=0mm] {\qihao\color{white}\bf {单一}};
	\draw (-0.35cm, 0.15 * 5.7cm) node[left, anchor = west, xshift=-1.5mm, yshift=0mm] {\bahao\color{gray2} {15}};
	\draw (0.4cm, 0.225 * 5.7cm ) node[left, anchor = west, yshift=0mm] {\qihao\color{white}\bf {较低}};
	\draw (-0.35cm, 0.3 * 5.7cm) node[left, anchor = west, xshift=-1.5mm, yshift=0mm] {\bahao\color{gray2} {30}};
	\draw (0.4cm, 0.65 * 5.7cm ) node[left, anchor = west, yshift=0mm] {\qihao\color{white}\bf {稳定}};
\end{tikzpicture}
\put(-2.5, 0){\raisebox{124.355 pt}{
\includegraphics[scale=2]{drop-l4.pdf}
}}
\put(11.5, 3.5){\raisebox{124.355 pt}{\qihao\color{gray2} {77}}}
\end{picture}}
}
\end{tctabularx}
\end{minipage}
\hfill
\noindent
\begin{minipage}[b]{.32\textwidth}
\fontsize{9.3pt}{10pt}\selectfont
\begin{tctabularx}{tabularx={p{4.6cm}}}
\parbox[c]{\hsize}{\centering\vskip4pt{\color{white}\bf 有益菌：
偏少 \raisebox{-0.3ex}{\includegraphics[scale=1]{redarrow_down.pdf}}
} \vskip4pt}
\\\\\\\\\\\\\\
\\\\\\\\\\\\\\
\\\\\\\\
\parbox[c]{\hsize}{
\hspace*{1cm}{
\begin{picture}(1, 1)
\begin{tikzpicture}[bar interval width=0.8, bar width=0.70cm]
	\begin{scope}
	\clip[] (0cm, 0cm) rectangle (1.5cm, 5.7cm);
	\draw[fill=solomon0, solomon0] (0cm, 0 * 5.7cm) rectangle (1.5cm, 0.05 * 5.7cm) node {};
	\draw[fill=solomon1, solomon1] (0cm, 0.05 * 5.7cm) rectangle (1.5cm, 0.2cm * 5.7cm) node {};
	\draw[fill=solomon3, solomon3] (0cm, 0.2 * 5.7cm) rectangle (1.5cm, 1 * 5.7cm) node {};
	\end{scope}
	\draw (-0.35cm, 0cm) node[left, anchor = west, yshift=0mm] {\bahao\color{gray2} {0}};
	\draw (0.4cm, 0.025 * 5.7cm ) node[left, anchor = west, yshift=0mm] {\qihao\color{white}\bf {缺乏}};
	\draw (-0.35cm, 0.05 * 5.7cm) node[left, anchor = west, yshift=0mm] {\bahao\color{gray2} {5}};
	\draw (0.4cm, 0.125 * 5.7cm ) node[left, anchor = west, yshift=0mm] {\qihao\color{white}\bf {偏少}};
	\draw (-0.35cm, 0.2 * 5.7cm) node[left, anchor = west, xshift=-1.5mm, yshift=0mm] {\bahao\color{gray2} {20}};
	\draw (0.4cm, 0.6 * 5.7cm ) node[left, anchor = west, yshift=0mm] {\qihao\color{white}\bf {健康}};
\end{tikzpicture}
\put(-2.5, 0){\raisebox{12.92 pt}{
\includegraphics[scale=2]{drop-l2.pdf}
}}
\put(13.5, 3.5){\raisebox{12.92 pt}{\qihao\color{gray2} {8}}}
\end{picture}}
}
\end{tctabularx}
\end{minipage}
\hfill
\noindent
\begin{minipage}[b][6cm]{.32\textwidth}
\fontsize{9.3pt}{10pt}\selectfont
\begin{tctabularx}{tabularx={p{4.6cm}}}
\parbox[c]{\hsize}{\centering\vskip4pt{\color{white}\bf 有害菌：
健康
} \vskip4pt}
\\\\\\\\\\\\\\
\\\\\\\\\\\\\\
\\\\\\\\
\parbox[c]{\hsize}{
\hspace*{1cm}{
\begin{picture}(1, 1)
\begin{tikzpicture}[bar interval width=0.8, bar width=0.70cm]
	\begin{scope}
	\clip[] (0cm, 0cm) rectangle (1.5cm, 5.7cm);
	\draw[fill=solomon3, solomon3] (0cm, 0 * 5.7cm) rectangle (1.5cm, 0.8 * 5.7cm) node {};
	\draw[fill=solomon1, solomon1] (0cm, 0.8 * 5.7cm) rectangle (1.5cm, 0.95cm * 5.7cm) node {};
	\draw[fill=solomon0, solomon0] (0cm, 0.95 * 5.7cm) rectangle (1.5cm, 1 * 5.7cm) node {};
	\end{scope}
	\draw (-0.35cm, 0cm) node[left, anchor = west, yshift=0mm] {\bahao\color{gray2} {0}};
	%\draw (1.45cm, 0cm ) node[left, anchor = west, yshift=0mm] {\includegraphics[scale=2]{drop-l1.pdf}};
	\draw (0.4cm, 0.4 * 5.7cm ) node[left, anchor = west, yshift=0mm] {\qihao\color{white}\bf {健康}};
	\draw (-0.35cm, 0.8 * 5.7cm) node[left, anchor = west, xshift=-1.5mm, yshift=0mm] {\bahao\color{gray2} {80}};
	\draw (0.4cm, 0.875 * 5.7cm ) node[left, anchor = west, yshift=0mm] {\qihao\color{white}\bf {偏高}};
	\draw (-0.35cm, 0.95 * 5.7cm) node[left, anchor = west, xshift=-1.5mm, yshift=0mm] {\bahao\color{gray2} {95}};
	\draw (0.5cm, 0.975 * 5.7cm ) node[left, anchor = west, yshift=0mm] {\qihao\color{white}\bf {高}};
\end{tikzpicture}
\put(-2.5, 0){\raisebox{1.615 pt}{
\includegraphics[scale=2]{drop-l4.pdf}
}}
\put(13.5, 3.5){\raisebox{1.615 pt}{\qihao\color{gray2} {1}}}
\end{picture}}
}
\end{tctabularx}
\end{minipage}

\vspace*{-0.5mm}
\hspace*{13cm}
\raisebox{-0.4ex}{\includegraphics{drop_mark.pdf}}
\fontsize{8.5pt}{11pt}\selectfont {本次检测结果}

%\vspace*{1mm}
%\noindent{\fontsize{11pt}{11pt}\selectfont\color{black70}\bf {本次检测异常结果}}

\vspace*{1.5mm}
\lantxh
\begin{tctabularxnew}{tabularx={m{4.84cm}<{\centering}m{4.84cm}<{\centering}m{4.84cm}<{\centering}}}
&&
\\[-6pt]
\cellcolor{topcolor} \raisebox{2.614pt}{\color{white}\bf 胖菌} &
\cellcolor{topcolor} \raisebox{2.614pt}{\color{white}\bf 瘦菌} &
\cellcolor{topcolor} \raisebox{2.614pt}{\color{white}\bf 致病菌}
\end{tctabularxnew}

\vspace*{-3.5mm}
\begin{spacing}{1.5}
\noindent\begin{tabular}{p{4.84cm}p{4.84cm}p{4.84cm}}
\hspace*{-2mm}\parbox[c]{1.08\hsize}{
\begin{allfour}
\fontsize{9pt}{11pt}\selectfont\color{black70}\centering {

\raisebox{-0.3ex}{\includegraphics[scale=1]{blank_arrow.pdf}}\\
未检出异常 \raisebox{-0.3ex}{\includegraphics[scale=1]{blank_arrow.pdf}}\\
\raisebox{-0.3ex}{\includegraphics[scale=1]{blank_arrow.pdf}}\\

}
\end{allfour}
} &
\hspace*{-1.5mm}\parbox[c]{1.08\hsize}{
\begin{allfour}
\fontsize{9pt}{11pt}\selectfont\color{black70}\centering {

\raisebox{-0.3ex}{\includegraphics[scale=1]{blank_arrow.pdf}}\\
未检出异常 \raisebox{-0.3ex}{\includegraphics[scale=1]{blank_arrow.pdf}}\\
\raisebox{-0.3ex}{\includegraphics[scale=1]{blank_arrow.pdf}}\\

}
\end{allfour}
} &
\hspace*{-1mm}\parbox[c]{1.08\hsize}{
\begin{allfour}
\fontsize{9pt}{11pt}\selectfont\color{black70}\centering {

\raisebox{-0.3ex}{\includegraphics[scale=1]{blank_arrow.pdf}}\\
未检出异常 \raisebox{-0.3ex}{\includegraphics[scale=1]{blank_arrow.pdf}}\\
\raisebox{-0.3ex}{\includegraphics[scale=1]{blank_arrow.pdf}}\\

}
\end{allfour}
}
\end{tabular}
\end{spacing}

\vspace*{-3mm}
\begin{spacing}{1.5}
\noindent\fontsize{8.5pt}{11pt}\selectfont {注：上表仅列出本次检测异常菌属。\raisebox{-0.3ex}{\includegraphics{black_up.pdf}}表示与参考人群比含量偏高/高；
\raisebox{-0.3ex}{\includegraphics{black_down.pdf}}表示与参考人群比含量偏低/低；致病菌仅列菌名表示检出但含量未超标。}
\end{spacing}

\vspace*{3mm}
\noindent{\fontsize{11pt}{11pt}\selectfont\color{black70}\bf {结果分析}}

\vspace*{0.5mm}
\begin{spacing}{1.5}
\noindent\color{black70}\fontsize{9pt}{12pt}\selectfont {{\raisebox{1pt}{\color{black70}$\bullet$}}\, 本次检测提示菌群多样性
正常、
有益菌
偏低、
有害菌
正常，
肠道菌群各项综合指标正常，整体状况良好。建议您保持良好的饮食生活习惯，规律作息，维持肠道菌群的稳态平衡。
}\\
\noindent\color{black70}\fontsize{9pt}{12pt}\selectfont {{\raisebox{1pt}{\color{black70}$\bullet$}}\,
本次检测未发现肠道菌群中胖菌及瘦菌相关菌属异常，胖菌/瘦菌通常与肠道健康、肥胖等健康问题相关。
}\\
\noindent\color{black70}\fontsize{9pt}{12pt}\selectfont {{\raisebox{1pt}{\color{black70}$\bullet$}}\,
本次未检测出致病菌。这是个好消息。
}
\end{spacing}
\end{document}

